\appendix

\section{Complete Prompt and Output Format}
\label{app:prompt}

All models received the same standardized answer--explanation instruction.
The required output structure was:

\begin{quote}
\texttt{FINAL\_ANSWER: <answer>}\\
\texttt{EXPLANATION: <concise justification>}
\end{quote}

The final answer was requested before the explanation to simplify answer
extraction and to support one common strict parser across all model--dataset
combinations.

The task-specific expected answer formats were:

\begin{itemize}
    \item \textbf{GSM8K}: a numerical answer or short numerical expression;
    \item \textbf{CommonsenseQA}: the selected option label, with option text
    permitted as additional content;
    \item \textbf{StrategyQA}: a binary yes/no answer.
\end{itemize}

The full prompt template used the following logical structure:

\begin{verbatim}
You are given a reasoning question.

Answer the question and provide a concise explanation.

Return exactly:
FINAL_ANSWER: <answer>
EXPLANATION: <concise justification>
\end{verbatim}

Dataset-specific question content and answer choices were inserted into this
shared template. No model-specific prompt engineering was used in the final
comparison.

\section{Models, Hardware, and Generation Configuration}
\label{app:configuration}

\subsection{Evaluated Models}

Table~\ref{tab:appendix_models} lists the exact model identifiers used in the
final experiments.

\begin{table*}[t]
\centering
\small
\caption{Evaluated open-weight language models.}
\label{tab:appendix_models}
\begin{tabular}{lll}
\toprule
\textbf{Short name} &
\textbf{Hugging Face identifier} &
\textbf{Precision} \\
\midrule
Qwen 2.5 3B &
\texttt{Qwen/Qwen2.5-3B-Instruct} &
4-bit NF4 \\

Llama 3.2 3B &
\texttt{meta-llama/Llama-3.2-3B-Instruct} &
4-bit NF4 \\

DeepSeek-R1 Distill 1.5B &
\texttt{deepseek-ai/DeepSeek-R1-Distill-Qwen-1.5B} &
4-bit NF4 \\
\bottomrule
\end{tabular}
\end{table*}

\subsection{Hardware and Software}

The experiments were conducted using:

\begin{itemize}
    \item Windows;
    \item Python~3.11.2;
    \item PyTorch~2.7.1 with CUDA~11.8;
    \item NVIDIA GeForce GTX~1650~Ti;
    \item 4\,GB VRAM.
\end{itemize}

Hugging Face Transformers and Datasets were used for model loading and data
processing. bitsandbytes and Accelerate were used for quantized model loading
and device placement. NumPy, Pandas, SciPy, and Matplotlib were used for
evaluation, statistical analysis, and visualization.

\subsection{Generation Settings}

Table~\ref{tab:appendix_generation} summarizes the final generation
configuration.

\begin{table}[t]
\centering
\small
\caption{Generation configuration used in the final experiments.}
\label{tab:appendix_generation}
\begin{tabular}{ll}
\toprule
\textbf{Parameter} & \textbf{Value} \\
\midrule
Decoding strategy & Greedy \\
Sampling & Disabled \\
Random seed & 42 \\
Qwen maximum new tokens & 256 \\
Llama maximum new tokens & 256 \\
DeepSeek maximum new tokens & 768 \\
Quantization & 4-bit NF4 \\
Double quantization & Enabled \\
Computation type & Float16 \\
\bottomrule
\end{tabular}
\end{table}

DeepSeek initially used the same 256-token limit as Qwen and Llama. Inspection
of raw outputs showed that many DeepSeek generations ended before both required
fields were completed. The final DeepSeek evaluation therefore used 768 new
tokens. This improved completion but reduced strict comparability of runtime
and output length.

\section{Manual Annotation Rubric and Failure Taxonomy}
\label{app:annotation}

\subsection{Sample Selection}

The reviewed manual-analysis sample contained 90 successfully parsed
responses:

\[
3\ \text{models}
\times
3\ \text{datasets}
\times
10\ \text{responses}
=
90\ \text{reviewed responses}.
\]

Ten responses were selected from each model--dataset combination. Five correct
and five incorrect answers were included whenever possible. Each model
therefore contributed 30 reviewed responses.

\subsection{Scoring Dimensions}

Each response was assessed on three dimensions:

\begin{itemize}
    \item \textbf{Explanation-support faithfulness}: whether the explanation
    correctly supports the model's stated final answer;

    \item \textbf{Logical consistency}: whether the reasoning is coherent and
    free from internal contradiction;

    \item \textbf{Completeness}: whether the explanation contains sufficient
    reasoning to justify the stated answer.
\end{itemize}

Each dimension was assigned a score from 0 to 3 according to
Table~\ref{tab:appendix_rubric}.

\begin{table}[t]
\centering
\small
\caption{Manual explanation-scoring rubric.}
\label{tab:appendix_rubric}
\begin{tabular}{cp{8.8cm}}
\toprule
\textbf{Score} & \textbf{Interpretation} \\
\midrule
0 &
The explanation is absent, unrelated, contradictory, or provides no
meaningful support for the final answer. \\

1 &
The explanation contains major factual, logical, or arithmetic errors, or
provides only weak support for the stated answer. \\

2 &
The explanation provides partial support but includes missing, unsupported,
incomplete, or flawed reasoning. \\

3 &
The explanation provides clear, correct, logically consistent, and sufficient
support for the final answer. \\
\bottomrule
\end{tabular}
\end{table}

\subsection{Primary Failure Categories}

A single primary failure category was assigned to each reviewed response:

\begin{description}
    \item[Answer--explanation mismatch:] the explanation supports a different
    answer from the reported final answer;

    \item[Arithmetic error:] the explanation contains an incorrect numerical
    operation or calculation;

    \item[Contradictory explanation:] different parts of the explanation are
    mutually inconsistent;

    \item[Factual hallucination:] the explanation introduces an unsupported or
    false factual claim;

    \item[Incomplete reasoning:] the explanation omits a necessary reasoning
    step or justification;

    \item[Instruction-following failure:] the response does not follow the
    required answer--explanation instruction;

    \item[Logical error:] the reasoning contains an invalid inference that is
    not primarily arithmetic;

    \item[Unsupported justification:] the explanation asserts a conclusion
    without sufficient support from the question or relevant knowledge;

    \item[Other:] a failure that does not fit the categories above;

    \item[No primary failure:] no major failure was identified.
\end{description}

The annotation procedure evaluates visible behavioral support. It does not
establish that the generated explanation is causally identical to the hidden
computation that produced the answer.

\section{Complete Accuracy Results}
\label{app:accuracy}

Table~\ref{tab:appendix_accuracy} reports answer accuracy with Wilson 95\%
confidence intervals for all model--dataset combinations.

\begin{table*}[t]
\centering
\small
\caption{Answer accuracy with Wilson 95\% confidence intervals.}
\label{tab:appendix_accuracy}
\begin{tabular}{llrrr}
\toprule
\textbf{Model} &
\textbf{Dataset} &
\textbf{Correct} &
\textbf{Accuracy (\%)} &
\textbf{95\% CI} \\
\midrule
Qwen 2.5 3B & GSM8K & 13/100 & 13.00 & [7.76, 20.98] \\
Qwen 2.5 3B & CommonsenseQA & 76/100 & 76.00 & [66.77, 83.31] \\
Qwen 2.5 3B & StrategyQA & 60/100 & 60.00 & [50.20, 69.06] \\
Qwen 2.5 3B & Overall & 149/300 & 49.67 & [44.05, 55.29] \\
\midrule
Llama 3.2 3B & GSM8K & 6/100 & 6.00 & [2.78, 12.48] \\
Llama 3.2 3B & CommonsenseQA & 70/100 & 70.00 & [60.42, 78.11] \\
Llama 3.2 3B & StrategyQA & 59/100 & 59.00 & [49.20, 68.13] \\
Llama 3.2 3B & Overall & 135/300 & 45.00 & [39.47, 50.66] \\
\midrule
DeepSeek-R1 Distill 1.5B & GSM8K & 80/100 & 80.00 & [71.12, 86.66] \\
DeepSeek-R1 Distill 1.5B & CommonsenseQA & 28/100 & 28.00 & [20.14, 37.49] \\
DeepSeek-R1 Distill 1.5B & StrategyQA & 49/100 & 49.00 & [39.42, 58.65] \\
DeepSeek-R1 Distill 1.5B & Overall & 157/300 & 52.33 & [46.69, 57.92] \\
\bottomrule
\end{tabular}
\end{table*}

The overall confidence intervals overlap substantially. Therefore, the
aggregate values should not be interpreted as decisive evidence that one model
is universally superior.

\section{Complete Faithfulness Statistics}
\label{app:statistics}

Table~\ref{tab:appendix_statistics} reports the comparison between
faithfulness scores for correct and incorrect answers.

\begin{table*}[t]
\centering
\small
\caption{Faithfulness comparison between correct and incorrect responses.}
\label{tab:appendix_statistics}
\begin{tabular}{lrrrrrr}
\toprule
\textbf{Scope} &
\textbf{Correct mean} &
\textbf{Incorrect mean} &
\textbf{Difference} &
\textbf{Bootstrap 95\% CI} &
\textbf{$p$-value} &
\textbf{Cliff's $\delta$} \\
\midrule
All models &
2.64 &
0.89 &
1.76 &
[1.42, 2.07] &
$<0.001$ &
0.804 \\

Qwen 2.5 3B &
2.73 &
0.73 &
2.00 &
[1.47, 2.47] &
$<0.001$ &
0.902 \\

Llama 3.2 3B &
2.67 &
0.73 &
1.93 &
[1.40, 2.40] &
$<0.001$ &
0.880 \\

DeepSeek-R1 Distill 1.5B &
2.53 &
1.20 &
1.33 &
[0.67, 1.93] &
0.001 &
0.627 \\
\bottomrule
\end{tabular}
\end{table*}

The corresponding Mann--Whitney U statistics were:

\begin{itemize}
    \item all models: $U=1827.0$;
    \item Qwen: $U=214.0$;
    \item Llama: $U=211.5$;
    \item DeepSeek: $U=183.0$.
\end{itemize}

All observed Cliff's delta values correspond to large effects. Positive values
indicate that explanations accompanying correct answers generally received
higher explanation-support faithfulness scores than explanations accompanying
incorrect answers.

\section{Additional Figures}
\label{app:additional_figures}

\begin{figure}[t]
    \centering
    \includegraphics[width=\linewidth]{faithfulness_heatmap.png}
    \caption{Mean explanation-support faithfulness by model and dataset.
    The values show substantial task-dependent variation.}
    \label{fig:appendix_heatmap}
\end{figure}

\begin{figure}[t]
    \centering
    \includegraphics[width=\linewidth]{failure_categories_by_model.png}
    \caption{Primary explanation failure categories in the balanced
    manual-analysis sample ($n=30$ per model).}
    \label{fig:appendix_failures}
\end{figure}

The heatmap shows that DeepSeek obtained the highest mean faithfulness on
CommonsenseQA (2.40) and GSM8K (2.10), but the lowest score on StrategyQA
(1.10). Llama obtained 2.00 on CommonsenseQA, 1.40 on GSM8K, and 1.70 on
StrategyQA. Qwen obtained 2.20, 1.30, and 1.70, respectively.

The failure-category figure shows that unsupported justification was common for
DeepSeek and Llama. Llama and Qwen also exhibited factual hallucinations and
answer--explanation mismatches, while Qwen additionally produced contradictory
explanations.

\section{Qualitative Examples}
\label{app:qualitative}

Table~\ref{tab:qualitative_examples} presents representative examples from the
reviewed manual-analysis sample. Excerpts are shortened for readability.

\begin{table*}[t]
\centering
\scriptsize
\caption{Representative qualitative examples from the reviewed sample.}
\label{tab:qualitative_examples}
\begin{tabular}{p{2.0cm}p{1.4cm}p{2.0cm}p{1.3cm}p{1.3cm}p{5.7cm}p{2.4cm}}
\toprule
\textbf{Model} &
\textbf{Dataset} &
\textbf{Case} &
\textbf{Gold} &
\textbf{Prediction} &
\textbf{Explanation excerpt} &
\textbf{Primary failure} \\
\midrule

DeepSeek-R1 Distill 1.5B &
CommonsenseQA &
Strong correct explanation &
E. tiredness &
E &
Tiredness is likely because running after a ball requires exerting energy,
leading to fatigue. &
None \\

Llama 3.2 3B &
GSM8K &
Answer--explanation mismatch &
720 &
960 &
Greta first puts 50\% of her \$2400 paycheck into retirement, which is
\$1200. She then uses 20\% of the paycheck for her car payment, which is
\$480. &
Answer--explanation mismatch \\

Qwen 2.5 3B &
CommonsenseQA &
Unsupported justification &
D. outside &
C. windowsill &
Ivy plants often require sunlight, which they can receive from warm
south- or west-facing windowsills. Windowsills also provide a stable
surface. &
Unsupported justification \\

Llama 3.2 3B &
StrategyQA &
Factual hallucination &
No &
B &
A modern CPU circuit chip is typically quite large and bulky, often
measuring several inches in size, whereas a house key is small and thin. &
Factual hallucination \\

DeepSeek-R1 Distill 1.5B &
GSM8K &
Logical or arithmetic error &
525 &
580 &
Andy needs to burn 105,000 calories over 181 days, which averages to
approximately 580 calories per day. &
Logical error \\
\bottomrule
\end{tabular}
\end{table*}

These examples illustrate that explanation fluency does not guarantee that the
reasoning is correct, complete, or aligned with the final answer.

\section{Reproducibility Details}
\label{app:reproducibility}

The final processed benchmark contained 300 questions and used random seed 42.
It was stored in both CSV and JSONL formats.

The SHA-256 checksum of the final processed benchmark CSV was:

\begin{quote}
\small
\texttt{35b721529e523d4a945407053ded18d93489e5476a23ec90fad1923fb36906e5}
\end{quote}

The complete experiment generated:

\[
3\ \text{models}
\times
3\ \text{datasets}
\times
100\ \text{samples}
=
900\ \text{responses}.
\]

Generated records preserved:

\begin{itemize}
    \item model identifier;
    \item dataset and sample identifier;
    \item complete prompt;
    \item raw generation;
    \item parsed final answer;
    \item extracted explanation;
    \item ground-truth answer;
    \item correctness label;
    \item generation time;
    \item strict parse status;
    \item processing error field.
\end{itemize}

Generated outputs were written incrementally to JSONL files so interrupted
runs could resume without regenerating completed samples. Model outputs were
evaluated separately by dataset and then merged into model-level and
cross-model result tables.

\section{Reproduction Commands}
\label{app:commands}

The main stages of the project were executed using commands of the following
form:

\begin{verbatim}
python -m src.prepare_data
python -m src.inspect_data
python -m src.run_inference
python -m src.evaluate_answers
python -m src.combine_results
python -m src.select_balanced_annotations
python -m src.analyze_manual_annotations
python -m src.statistical_analysis
python -m src.select_qualitative_examples
python -m src.create_paper_figures
\end{verbatim}

The exact model, dataset, input, output, token-budget, and overwrite arguments
are preserved in the project scripts and configuration files.
